% sec3_6.tex — Section 3.6: Testing Overidentifying Restrictions

If the equation is exactly identified, then it is possible to choose $\tilde{\boldsymbol{\delta}}$
so that all the elements of the sample moments $\mathbf{g}_n(\tilde{\boldsymbol{\delta}})$ are zero and
the distance
\[
  J(\tilde{\boldsymbol{\delta}},\, \hat{\mathbf{W}})
  \equiv n \cdot \mathbf{g}_n(\tilde{\boldsymbol{\delta}})'\,
  \hat{\mathbf{W}}\, \mathbf{g}_n(\tilde{\boldsymbol{\delta}})
\]
is zero.  (The $\tilde{\boldsymbol{\delta}}$ that does it is the IV estimator.)  If the equation is
overidentified, then the distance cannot be set to zero exactly, but we would expect the
minimized distance to be close to zero.  It turns out that, if the weighting matrix
$\hat{\mathbf{W}}$ is chosen optimally so that $\plim\, \hat{\mathbf{W}} = \mathbf{S}^{-1}$, then
the minimized distance is asymptotically chi-squared.

Let $\hat{\mathbf{S}}$ be a consistent estimator of $\mathbf{S}$, and consider first the case where the
distance is evaluated at the true parameter value $\boldsymbol{\delta}$,
$J(\boldsymbol{\delta}, \hat{\mathbf{S}}^{-1})$.  Since by definition
$\mathbf{g}_n(\boldsymbol{\delta}) = \bar{\mathbf{g}}$
$(\equiv \frac{1}{n} \sum_i \mathbf{g}_i)$ for
$\tilde{\boldsymbol{\delta}} = \boldsymbol{\delta}$, the distance equals
\begin{equation}
  J(\boldsymbol{\delta}, \hat{\mathbf{S}}^{-1})
  = n \cdot \bar{\mathbf{g}}'\, \hat{\mathbf{S}}^{-1}\, \bar{\mathbf{g}}
  = (\sqrt{n}\, \bar{\mathbf{g}})'\, \hat{\mathbf{S}}^{-1}\,
    (\sqrt{n}\, \bar{\mathbf{g}}).
  \label{eq:3.6.1}
\end{equation}
Since $\sqrt{n}\, \bar{\mathbf{g}} \dto N(\mathbf{0}, \mathbf{S})$ and
$\hat{\mathbf{S}} \pto \mathbf{S}$, its asymptotic distribution is $\chi^2(K)$ by
Lemma~2.4(d).  Now if $\boldsymbol{\delta}$ is replaced by
$\hat{\boldsymbol{\delta}}(\hat{\mathbf{S}}^{-1})$, then the degrees of freedom change
from $K$ to $K - L$.  The intuitive reason is that we have to estimate $L$ parameters
$\boldsymbol{\delta}$ before forming the sample average of $\mathbf{g}_i$.  (We encountered a similar
situation in Chapter~1 in the context of the unbiased estimation of $\sigma^2$.)  We summarize
this as

\begin{proposition}[Hansen's test of overidentifying restrictions (Hansen, 1982)]
\label{prop:3.6}
\emph{Suppose there is available a consistent estimator, $\hat{\mathbf{S}}$, of
$\mathbf{S}$ $(= \E(\mathbf{g}_i \mathbf{g}_i'))$.  Under Assumptions~3.1--3.5,}
\[
  J\bigl(\hat{\boldsymbol{\delta}}(\hat{\mathbf{S}}^{-1}),\,
  \hat{\mathbf{S}}^{-1}\bigr)
  \;\bigl(= n \cdot
  \mathbf{g}_n(\hat{\boldsymbol{\delta}}(\hat{\mathbf{S}}^{-1}))'\,
  \hat{\mathbf{S}}^{-1}\,
  \mathbf{g}_n(\hat{\boldsymbol{\delta}}(\hat{\mathbf{S}}^{-1}))\bigr)
  \dto \chi^2(K - L).
\]
\end{proposition}

A formal proof is left as an optional exercise.  Because the $\hat{\mathbf{S}}$ given in
\eqref{eq:3.5.10} is consistent (under the additional condition of Assumption~3.6), the
minimized distance calculated in the second step of the efficient two-step GMM is
asymptotically $\chi^2(K - L)$.

Three points on the use and interpretation of the $J$ test are in order.

\begin{itemize}
\item This is a \textbf{specification test}, testing whether all the restrictions of the model
  (which are the assumptions maintained in Proposition~3.6) are satisfied.  If the $J$
  statistic of Proposition~3.6 is surprisingly large, it means that either the orthogonality
  conditions (Assumption~3.3) or the other assumptions (or both) are likely
  to be false.  Only when we are confident about those other assumptions can we
  interpret the large $J$ statistic as evidence for the endogeneity of some of the $K$
  instruments included in $\mathbf{x}_i$.

\item Unlike the tests we have encountered so far, the test is not consistent against
  some failures of the orthogonality conditions.  The essential reason is the loss of
  degrees of freedom from $K$ to $K - L$.  It is easy to show that $\bar{\mathbf{g}}$ is related to its
  sample counterpart, $\mathbf{g}_n(\hat{\boldsymbol{\delta}}(\hat{\mathbf{S}}^{-1}))$, as
  \begin{equation}
    \sqrt{n}\, \mathbf{g}_n(\hat{\boldsymbol{\delta}}(\hat{\mathbf{S}}^{-1}))
    = \hat{\mathbf{B}}\, \sqrt{n}\, \bar{\mathbf{g}}, \quad
    \hat{\mathbf{B}} \equiv \mathbf{I}_K
    - \mathbf{S}_{\mathbf{xz}}
      (\mathbf{S}_{\mathbf{xz}}' \hat{\mathbf{S}}^{-1} \mathbf{S}_{\mathbf{xz}})^{-1}
      \mathbf{S}_{\mathbf{xz}}' \hat{\mathbf{S}}^{-1}.
    \label{eq:3.6.2}
  \end{equation}
  The problem is that, since
  $\hat{\mathbf{B}} \mathbf{S}_{\mathbf{xz}} = \mathbf{0}$, this matrix $\hat{\mathbf{B}}$ is not of full
  column rank.  If the orthogonality conditions fail and $\E(\mathbf{g}_i) \neq \mathbf{0}$, then the
  elements of $\sqrt{n}\, \bar{\mathbf{g}}$ will diverge to $+\infty$ or $-\infty$.  But, since
  $\hat{\mathbf{B}}$ is not of full column rank,
  $\hat{\mathbf{B}} \sqrt{n}\, \bar{\mathbf{g}}$ and hence
  $J(\hat{\boldsymbol{\delta}}(\hat{\mathbf{S}}^{-1}), \hat{\mathbf{S}}^{-1})$ may remain finite
  for some pattern of nonzero orthogonalities.\footnote{See Newey (1985, Section~3) for a
  thorough treatment of this issue.}

\item It is only recently that the small-sample properties of the test became a concern.
  Several papers in the July 1996 issue of the \emph{Journal of Business and Economic
  Statistics} report that the finite-sample or actual size of the $J$ test in small samples
  far exceeds the nominal size (i.e., the test rejects too often).
\end{itemize}

%% ─── Testing Subsets of Orthogonality Conditions ───────────────────────────
\subsection*{Testing Subsets of Orthogonality Conditions}

Suppose we can divide the $K$ instruments into two groups: the vector $\mathbf{x}_{i1}$ of $K_1$
variables that are known to satisfy the orthogonality conditions, and the vector
$\mathbf{x}_{i2}$ of remaining $K - K_1$ variables that are suspect.  Since the ordering of instruments
does not change the numerical values of the estimator and test statistics, we can
assume without loss of generality that the last $K - K_1$ elements of $\mathbf{x}_i$ are the suspect
instruments:
\begin{equation}
  \mathbf{x}_i = \begin{bmatrix} \mathbf{x}_{i1} \\ \mathbf{x}_{i2} \end{bmatrix}
  \quad
  \left\}
  \begin{array}{l} K_1 \text{ rows,} \\ K - K_1 \text{ rows.} \end{array}
  \right.
  \label{eq:3.6.3}
\end{equation}
The part of the model we wish to test is
\begin{equation}
  \E(\mathbf{x}_{i2} \cdot \varepsilon_i) = \mathbf{0}.
  \label{eq:3.6.4}
\end{equation}
This restriction is testable if there are at least as many nonsuspect instruments as
there are coefficients so that $K_1 \geq L$.  The basic idea is to compare two $J$ statistics
from two separate GMM estimators of the same coefficient vector $\boldsymbol{\delta}$, one using only
the instruments included in $\mathbf{x}_{i1}$, and the other using also the suspect instruments
$\mathbf{x}_{i2}$ in addition to $\mathbf{x}_{i1}$.  If the inclusion of the suspect instruments significantly
increases the $J$ statistic, that is a good reason for doubting the predeterminedness of
$\mathbf{x}_{i2}$.

In accordance with the partition of $\mathbf{x}_i$, the sample orthogonality conditions
$\mathbf{g}_n(\tilde{\boldsymbol{\delta}})$ and $\mathbf{S}$ can be written as
\begin{equation}
  \underset{(K \times 1)}{\mathbf{g}_n(\tilde{\boldsymbol{\delta}})}
  = \begin{bmatrix}
      \mathbf{g}_{1n}(\tilde{\boldsymbol{\delta}}) \\[-2pt]
      {\scriptstyle (K_1 \times 1)} \\[4pt]
      \mathbf{g}_{2n}(\tilde{\boldsymbol{\delta}}) \\[-2pt]
      {\scriptstyle ((K - K_1) \times 1)}
    \end{bmatrix}, \quad
  \underset{(K \times K)}{\mathbf{S}}
  = \begin{bmatrix}
      \mathbf{S}_{11} & \mathbf{S}_{12} \\
      \mathbf{S}_{21} & \mathbf{S}_{22}
    \end{bmatrix},
  \label{eq:3.6.5}
\end{equation}
where
\[
  \mathbf{S}_{11} = \E(\varepsilon_i^2\, \mathbf{x}_{i1} \mathbf{x}_{i1}'), \quad
  \mathbf{S}_{12} = \E(\varepsilon_i^2\, \mathbf{x}_{i1} \mathbf{x}_{i2}'), \quad
  \mathbf{S}_{21} = \E(\varepsilon_i^2\, \mathbf{x}_{i2} \mathbf{x}_{i1}'), \quad
  \mathbf{S}_{22} = \E(\varepsilon_i^2\, \mathbf{x}_{i2} \mathbf{x}_{i2}').
\]
In particular, $\mathbf{g}_{1n}(\tilde{\boldsymbol{\delta}})$ can be written as
\[
  \mathbf{g}_{1n}(\tilde{\boldsymbol{\delta}})
  = \frac{1}{n} \sum_{i=1}^{n} \mathbf{x}_{i1} \cdot
    (y_i - \mathbf{z}_i' \tilde{\boldsymbol{\delta}})
  \equiv \mathbf{s}_{\mathbf{x}_1 \mathbf{y}}
  - \mathbf{S}_{\mathbf{x}_1 \mathbf{z}}\, \tilde{\boldsymbol{\delta}},
\]
where
\begin{equation}
  \mathbf{s}_{\mathbf{x}_1 \mathbf{y}}
  \equiv \frac{1}{n} \sum_{i=1}^{n} \mathbf{x}_{i1} \cdot y_i, \quad
  \mathbf{S}_{\mathbf{x}_1 \mathbf{z}}
  \equiv \frac{1}{n} \sum_{i=1}^{n} \mathbf{x}_{i1}\, \mathbf{z}_i'.
  \label{eq:3.6.6}
\end{equation}

For a consistent estimate $\hat{\mathbf{S}}$ of $\mathbf{S}$, the efficient GMM estimator using all the $K$
instruments and its associated $J$ statistic have already been derived in this and the
previous section.  Reproducing them,
\begin{align}
  \hat{\boldsymbol{\delta}}
  &= (\mathbf{S}_{\mathbf{xz}}' \hat{\mathbf{S}}^{-1} \mathbf{S}_{\mathbf{xz}})^{-1}\,
     \mathbf{S}_{\mathbf{xz}}' \hat{\mathbf{S}}^{-1} \mathbf{s}_{\mathbf{xy}},
  \label{eq:3.6.7} \\
  J &= n \cdot \mathbf{g}_n(\hat{\boldsymbol{\delta}})'\,
     \hat{\mathbf{S}}^{-1}\, \mathbf{g}_n(\hat{\boldsymbol{\delta}}).
  \label{eq:3.6.8}
\end{align}
The efficient GMM estimator of the \emph{same} coefficient vector $\boldsymbol{\delta}$ using only the first
$K_1$ instruments and its associated $J$ statistic are obtained by replacing $\mathbf{x}_i$ by
$\mathbf{x}_{i1}$ in these expressions.  So
\begin{align}
  \tilde{\boldsymbol{\delta}}
  &= [\mathbf{S}_{\mathbf{x}_1 \mathbf{z}}'
     (\hat{\mathbf{S}}_{11})^{-1}
     \mathbf{S}_{\mathbf{x}_1 \mathbf{z}}]^{-1}\,
     \mathbf{S}_{\mathbf{x}_1 \mathbf{z}}'
     (\hat{\mathbf{S}}_{11})^{-1}
     \mathbf{s}_{\mathbf{x}_1 \mathbf{y}},
  \label{eq:3.6.9} \\
  J_1 &= n \cdot \mathbf{g}_{1n}(\tilde{\boldsymbol{\delta}})'\,
     (\hat{\mathbf{S}}_{11})^{-1}\,
     \mathbf{g}_{1n}(\tilde{\boldsymbol{\delta}}),
  \label{eq:3.6.10}
\end{align}
where $\hat{\mathbf{S}}_{11}$ is a consistent estimate of $\mathbf{S}_{11}$.

The test is based on the following proposition specifying the asymptotic distribution
of $J - J_1$ (the proof is left as an optional exercise).

\begin{proposition}[testing a subset of orthogonality conditions\footnote{The test was
developed in Newey (1985, Section~4) and Eichenbaum, Hansen, and Singleton (1985,
Appendix~C).}]
\label{prop:3.7}
\emph{Suppose Assumptions~3.1--3.5 hold.  Let $\mathbf{x}_{i1}$ be a subvector of $\mathbf{x}_i$, and
strengthen Assumption~3.4 by requiring that the rank condition for identification is
satisfied for $\mathbf{x}_{i1}$ (so $\E(\mathbf{x}_{i1} \mathbf{z}_i')$ is of full column rank).  Then,
for any consistent estimators $\hat{\mathbf{S}}$ of $\mathbf{S}$ and $\hat{\mathbf{S}}_{11}$ of
$\mathbf{S}_{11}$,}
\[
  C \equiv J - J_1 \dto \chi^2(K - K_1),
\]
\emph{where $K = \#\mathbf{x}_i$ (dimension of $\mathbf{x}_i$),
$K_1 = \#\mathbf{x}_{i1}$ (dimension of $\mathbf{x}_{i1}$), and $J$ and $J_1$ are defined in
\eqref{eq:3.6.8} and \eqref{eq:3.6.10}.}
\end{proposition}

Clearly, the choice of $\hat{\mathbf{S}}$ and $\hat{\mathbf{S}}_{11}$ does not matter asymptotically as long
as they are consistent.  But in finite samples, the test statistic $C$ can be negative.  This
problem can be avoided and $C$ can be made nonnegative in finite samples if the same
$\hat{\mathbf{S}}$ is used throughout, that is, if $\hat{\mathbf{S}}_{11}$ in \eqref{eq:3.6.9} and
\eqref{eq:3.6.10} is the submatrix of $\hat{\mathbf{S}}$ in \eqref{eq:3.6.7} and \eqref{eq:3.6.8}.
This is accomplished by taking the following steps:

\begin{enumerate}
\item[(1)] do the efficient two-step GMM with full instruments $\mathbf{x}_i$ to obtain
  $\hat{\mathbf{S}}$ from the first step, $\hat{\boldsymbol{\delta}}$ and $J$ from the second step;

\item[(2)] extract the leading $K_1 \times K_1$ submatrix $\hat{\mathbf{S}}_{11}$ from $\hat{\mathbf{S}}$
  obtained from (1), calculate $\tilde{\boldsymbol{\delta}}$ by \eqref{eq:3.6.9} using this
  $\hat{\mathbf{S}}_{11}$ and $J_1$ by \eqref{eq:3.6.10}.  Then take the difference in $J$.
\end{enumerate}

It is left as an optional analytical exercise to prove that $C$ calculated as just described
is nonnegative in finite samples.

We can use Proposition~3.7 to test for the endogeneity of a subset of regressors,
as the next example illustrates.

\begin{example}[testing whether schooling is predetermined in the wage equation]
\label{ex:3.3}
In the wage equation of Example~3.2, suppose schooling $S_i$ is
suspected to be endogenous.  To test for the endogeneity of $S_i$, partition $\mathbf{x}_i$ as
\[
  \mathbf{x}_{i1} = \begin{bmatrix} 1 \\ \mathit{EXPR}_i \\ \mathit{AGE}_i \\
  \mathit{MED}_i \end{bmatrix}, \quad
  \mathbf{x}_{i2} = S_i.
\]
The vector of regressors, $\mathbf{z}_i$, is the same as in Example~3.2.  The first step is to
do the efficient two-step GMM estimation of $\boldsymbol{\delta}$ with
$\mathbf{x}_i \equiv (1, \mathit{EXPR}_i, \mathit{AGE}_i, \mathit{MED}_i, S_i)'$ as the instruments.
This produces $J$ and the $5 \times 5$ matrix $\hat{\mathbf{S}}$.
Second, extract the leading $4 \times 4$ submatrix $\hat{\mathbf{S}}_{11}$ corresponding to
$\mathbf{x}_{i1}$ from $\hat{\mathbf{S}}$ and estimate the same coefficients $\boldsymbol{\delta}$ by GMM, this
time with the fewer instruments $\mathbf{x}_{i1}$ and using this $\hat{\mathbf{S}}_{11}$.  The difference
in the $J$ statistic from the two different GMM estimators of $\boldsymbol{\delta}$ should be
asymptotically $\chi^2(1)$.
\end{example}

%% ─── Questions for Review ──────────────────────────────────────────────────
\bigskip
\begin{center}\rule{0.6\textwidth}{0.4pt}\end{center}
\noindent\textsc{Questions for Review}
\medskip

\begin{enumerate}
\item Does
  $J(\hat{\boldsymbol{\delta}}(\hat{\mathbf{S}}^{-1}),\, \bar{\mathbf{S}}^{-1})
  \dto \chi^2(K - L)$ where $\hat{\mathbf{S}}$ and $\bar{\mathbf{S}}$ are two different consistent
  estimators of $\mathbf{S}$?

\item Show:
  $J(\hat{\boldsymbol{\delta}}(\hat{\mathbf{S}}^{-1}),\, \hat{\mathbf{S}}^{-1})
  = n \cdot \mathbf{s}_{\mathbf{xy}}'\, \hat{\mathbf{S}}^{-1}\,
  (\mathbf{s}_{\mathbf{xy}} - \mathbf{S}_{\mathbf{xz}}\,
  \hat{\boldsymbol{\delta}}(\hat{\mathbf{S}}^{-1}))$.

\item Can the degrees of freedom of the $C$ statistic be greater than $K - L$?
  [Answer: No.]

\item Suppose $K_1 = L$.  Does the numerical value of $C$ depend on the partition of
  $\mathbf{x}_i$ between $\mathbf{x}_{i1}$ and $\mathbf{x}_{i2}$?
\end{enumerate}
